% Abhakseite „Das kann ich“ – Zone nur im Gesamt (\mitzone)
%% TODO Ich-kann-Sätze: Platzhalter wie in den Aufgabendateien
\begin{abhakseite}
\ifdefined\mitzone
\abhakgruppe{Kennst du schon}
\abhak{1}{Bruch als Anteil lesen, am Streifen einzeichnen (drei Achtel markieren)}
\abhak{2}{Kürzen und Erweitern; gleichwertige Brüche (zwei Drittel gleich vier Sechstel)}
\abhak{3}{Vielfache und Teiler (kgV von vier und sechs, ggT von zwölf und achtzehn)}
\abhak{4}{Stellenwerte der Dezimalzahlen (Zehntel, Hundertstel; 0,7 = 0,70)}
\abhak{5}{Schriftliches Rechnen mit natürlichen Zahlen, Einmaleins}
\abhak{6}{Bruch ↔ Dezimalzahl bei einfachen Brüchen (ein Halb, drei Viertel)}
\abhak{7}{Kürzen und Erweitern; gleichwertige Brüche (zwei Drittel gleich vier Sechstel) – Fehler finden}
\abhak{8}{Kürzen und Erweitern; gleichwertige Brüche (zwei Drittel gleich vier Sechstel) – selbst rechnen}
\fi
\abhakgruppe{Einheit 1 · Brüche addieren und subtrahieren}
\abhak{9}{Gleicher Nenner – ja oder nein?}
\abhak{10}{Womit erweitere ich?}
\abhak{11}{Addieren/Subtrahieren}
\abhak{12}{ganze Zahl plus Bruch}
\abhak{13}{Addieren/Subtrahieren – Fehler finden, Begründen, Anwendung, Darstellungswechsel}
\abhakgruppe{Einheit 2 · Dezimalzahlen addieren und subtrahieren}
\abhak{14}{Dezimal addieren/subtrahieren}
\abhak{15}{ganze Zahl plus Dezimalzahl}
\abhak{16}{Überschlag vorab}
\abhak{17}{Dezimal addieren/subtrahieren – Fehler finden, Begründen, Anwendung}
\abhakgruppe{Einheit 3 · Brüche multiplizieren und dividieren}
\abhak{18}{Von heißt mal}
\abhak{19}{Multiplizieren}
\abhak{20}{Dividieren}
\abhak{21}{Multiplizieren – Fehler finden, Anwendung, Darstellungswechsel}
\abhak{22}{Dividieren – Fehler finden, Begründen}
\abhakgruppe{Einheit 4 · Dezimalzahlen multiplizieren und dividieren}
\abhak{23}{Wie viele Kommastellen hat das Ergebnis?}
\abhak{24}{Dezimal multiplizieren/dividieren}
\abhak{25}{Dezimal multiplizieren/dividieren – Fehler finden, Begründen, Anwendung}
\abhakgruppe{Einheit 5 · Rechengesetze und Punkt vor Strich}
\abhak{26}{Punkt vor Strich}
\abhak{27}{Distributivgesetz (Ausklammern bei Dezimalzahlen)}
\abhak{28}{Überschlag und Prüfen}
\abhak{29}{Punkt vor Strich – Fehler finden, Begründen, Anwendung}
\end{abhakseite}
